\subsubsection{The traditional Form of \bsq{switch}}
Originally, a \lstinline|switch|\index{switch}\index{switch!statement} statement in \Java looked like this (the ``\verb#case L:#'' variant):
\begin{lstlisting}[numbers=left]
switch( <selector> )
{
    case <switchlabel1>:
        <statements1>;
        // Falls through!

    case <switchlabel2>:
        <statements2>;
        break;

    case <switchlabel3>: <singleStatement>; break;

    case <switchlabel4>:
    {
        <statements>;
        break;
    }

    case <switchlabel5>: // Also a fall-through
    case <switchlabel6>:
    {
        <statements>;
        break;
    }

    default:
       <statements>;
       break;
}
\end{lstlisting}

First, these coding conventions make the blank line above each \lstinline|case|\index{case} (except the very first one) and above \lstinline|default| mandatory!
  
If a \lstinline|case| falls through, like in lines~3 to 9, a comment \bdq{\lstinline|// Falls through!|} (or something similar\footnote{Linting tools as well as most IDEs issue a warning in case of a fall-through. Usually, these could be switched off by adding a comment or an annotation to the location of the fall-through. This would be a sufficient replacement for the comment.}) as in line~5 of the sample is required! But basically, such constructs should be avoided when possible, because it also forces that the sequence of the \lstinline|case| clauses cannot be changed, and that fact requires an additional comment to the \lstinline|switch| itself. In fact here you should consider to ignore the \bdq{DRY~Principle}\index{DRY Principle} – we will see this later – and repeat the code for the second \lstinline|case| on the first one:
\begin{lstlisting}
// BETTER
switch( <selector> )
{
    case <switchlabel1>:
        <statements1>;
        <statements2>;
        break;

    case <switchlabel2>:
        <statements2>;
        break;

    default: …
}
\end{lstlisting}
Or you put the code that is common for both cases into a \lstinline|private| method and call that from both branches. 

Different to the case shown in lines~3 to 9, no comment is required in the case shown in lines~19 and 20: here we have two \lstinline|case| clauses that triggers \textit{exactly the same} action, and we can even omit the blank line between them (and also the comment from line~19).

For longer \lstinline|switch| statements it is highly recommended to use a label with the \lstinline|break| statement\footnote{In fact, I recommend to always use a label with \lstinline|break|; see chapter \tqfullref{sec:LabelsAndBreakStatements} on this topic.}; this can look like this:
\begin{lstlisting}
final Color color;
ColorSelector: switch( colorIndex )
{
    case 1: color = RED; 
        break ColorSelector;
        
    case 2: color = BLUE; 
        break ColorSelector;
        
    case 3: color = YELLOW; 
        break ColorSelector;
        
    case 4: color = GREEN; 
        break ColorSelector;
        
    default: color = NONE; 
        break ColorSelector;
}   //  ColorSelector:
\end{lstlisting}

The \lstinline|break| in the \lstinline|default| branch is redundant, but it prevents a fall-through error if later another \lstinline|case| branch is added – although it should be assured that the \lstinline|default| branch is always the last branch in the \lstinline|switch| statement. Of course no \lstinline|break| is required when the branch is left by throwing an exception.\footnote{A \lstinline|break| would be also obsolete when the branch is left through a \lstinline|return| statement, but in that case, the \lstinline|return| would not be the last statement in the method, and this is another requirement from this coding conventions!}

%==============================================================================
% End of file!!
%==============================================================================
